\documentclass[11pt,a4paper]{article}

% --- Paquetes Esenciales ---
\usepackage[utf8]{inputenc}
\usepackage[T1]{fontenc}
\usepackage{lmodern}
\usepackage{microtype}

% Renombrar elementos a español sin depender de babel-spanish
\renewcommand{\contentsname}{Índice General}
\renewcommand{\abstractname}{Resumen}
\renewcommand{\tablename}{Tabla}
\renewcommand{\refname}{Referencias}

% --- Geometría y Márgenes ---
\usepackage[top=2.5cm, bottom=2.5cm, left=2.5cm, right=2.5cm]{geometry}

% --- Colores y Gráficos ---
\usepackage{xcolor}
\usepackage{graphicx}
\usepackage{tikz}

% Definición de Paleta Profesional
\definecolor{primary}{RGB}{24, 43, 73}      % Azul marino institucional
\definecolor{secondary}{RGB}{41, 128, 185}  % Azul cobalto
\definecolor{darkgray}{RGB}{55, 65, 81}
\definecolor{lightgray}{RGB}{243, 244, 246}
\definecolor{codebg}{RGB}{249, 250, 251}
\definecolor{accent}{RGB}{217, 83, 79}

% --- Matemáticas ---
\usepackage{amsmath, amssymb, amsfonts, amsthm}
\usepackage{mathtools}

% --- Tablas ---
\usepackage{booktabs}
\usepackage{tabularx}
\usepackage{array}
\usepackage{multirow}

% --- Cajas y Énfasis (tcolorbox) ---
\usepackage[most]{tcolorbox}
\newtcolorbox{infobox}[1][]{
  colback=lightgray,
  colframe=primary,
  fonttitle=\bfseries\color{white},
  coltitle=white,
  boxrule=0.8pt,
  arc=2.5mm,
  left=4mm, right=4mm, top=3mm, bottom=3mm,
  title=#1
}

% --- Encabezados y Pies de Página ---
\usepackage{fancyhdr}
\pagestyle{fancy}
\fancyhf{}
\fancyhead[L]{\small\color{darkgray}\textit{Modelado y Programacion}}
\fancyhead[R]{\small\color{darkgray}\thepage}
\fancyfoot[L]{\small\color{darkgray}\textsf{Facultad de Ciencias}}
\fancyfoot[R]{\small\color{darkgray}\textsf{Septiembre 2026}}
\renewcommand{\headrulewidth}{0.4pt}
\renewcommand{\footrulewidth}{0.4pt}

% --- Títulos de Secciones (titlesec) ---
\usepackage{titlesec}
\titleformat{\section}
  {\color{primary}\normalfont\Large\bfseries}
  {\color{primary}\thesection}{1em}{}[\color{secondary}\titlerule]
\titleformat{\subsection}
  {\color{secondary}\normalfont\large\bfseries}
  {\thesubsection}{1em}{}
\titleformat{\subsubsection}
  {\color{darkgray}\normalfont\normalsize\bfseries}
  {\thesubsubsection}{1em}{}

% --- Código Fuente (listings) ---
\usepackage{listings}
\lstset{
  backgroundcolor=\color{codebg},
  basicstyle=\ttfamily\footnotesize\color{darkgray},
  keywordstyle=\color{secondary}\bfseries,
  commentstyle=\color{gray}\itshape,
  stringstyle=\color{accent},
  numbers=left,
  numberstyle=\tiny\color{gray},
  stepnumber=1,
  numbersep=8pt,
  frame=single,
  rulecolor=\color{lightgray!80!black},
  breaklines=true,
  showstringspaces=false,
  captionpos=b
}

% --- Enlaces e Hipervínculos ---
\usepackage[colorlinks=true,
            linkcolor=primary,
            citecolor=secondary,
            urlcolor=secondary]{hyperref}

% --- Metadatos del Documento ---
\newcommand{\reporttitle}{Análisis de Requisitos Funcionales del Proyecto}
\newcommand{\reportsubtitle}{Documento de Especificaciones}

\newcommand{\reportauthor}{
  -Bárcenas Lejarazo Karen \\[0.1cm]

  -Cantero Zavaleta Héctor \\[0.1cm]

  -Miranda Vieyra Francisco Cuahutémoc \\[0.1cm]

  -Zúñiga Vilchis Carlos Uriel
}

\newcommand{\reportinstitution}{Facultad de Ciencias / UNAM}
\newcommand{\reportcourse}{Modelado y Programacion}
\newcommand{\reportdate}{\today}

\begin{document}

\begin{titlepage}
    \centering
    \vspace*{1.2cm}
    
    % Barra de acento visual superior
    {\color{secondary}\rule{\textwidth}{2.5pt}}\\[0.8cm]
    
    {\Huge \textbf{\color{primary}\reporttitle}}\\[0.5cm]
    {\Large \textit{\color{darkgray}\reportsubtitle}}\\[0.8cm]
    
    {\color{secondary}\rule{\textwidth}{1pt}}\\[2.2cm]
    
    \begin{minipage}{0.48\textwidth}
        \begin{flushleft} \large
            \textbf{\color{primary}Autor(es):}\\
            \reportauthor\\[0.2cm]
        \end{flushleft}
    \end{minipage}
    \hfill
    \begin{minipage}{0.48\textwidth}
        \begin{flushright} \large
            \textbf{\color{primary}Institución / Asignatura:}\\
            \reportinstitution\\
            \reportcourse\\[0.2cm]
            \textbf{\color{primary}Docente / Asesor:}\\
            José de Jesús Galaviz Casas
        \end{flushright}
    \end{minipage}
    
    \vfill
    
    \begin{tcolorbox}[colback=lightgray!60, colframe=lightgray!80!black, arc=2mm, width=0.88\textwidth, halign=center]
        \textbf{Fecha de Publicación / Entrega:} \reportdate \quad $\vert$ \quad \textbf{Versión:} 1.0
    \end{tcolorbox}
    
    \vspace{0.8cm}
\end{titlepage}

% ==========================================================
% SECCIÓN 1: Introduccion
% ==========================================================

\section{Introducción}
Para poder realizar una correcta construccion de cualquier aplicación software es necesario poder conocer de manera precisa el problema al cual nos
enfrentamos, ya que de esta manera se espera que la solución se desarrolle bajo ese conocimiento sea las mas adecuada y útil. Debido a esto, es necesario realizar un correcto analisis de los requerimientos del sistema, puesto que esto determinara lo que hara nuestro programa, las limitaciones tecnicas y la forma de operacion de este.

\subsection{Objetivos del Proyecto}
\begin{itemize}
    \item \textbf{Objetivo Primario:} Localizar de manera correcta los requisitos funcionales y no funcionales del proyecto.
    \item \textbf{Objetivos Específicos:}
    \begin{enumerate}
        \item Realizar una separacion correcta entre la funcionalidad y la no funcionalidad.
        \item Lograr un analisis detallado del proyecto para captar los tipos de requisitos.
        \item Contrastar las diferencias entre cada requisito.
    \end{enumerate}
\end{itemize}

\section{Los distintos requisitos funcionales del proyecto}




\end{document}

